\section{One Radio, Three Protocols}
\label{sec:ble-one-radio}

Everything in this chapter so far treats the radio as a Bluetooth radio. It is not. It is a 2.4\,GHz radio
that happens to be running Bluetooth at the moment, and the same silicon is meant to run 802.15.4 --- Zigbee
and Thread --- later. That is not an aspiration, it is measurable in the archive this project already
links:

\begin{minted}[fontsize=\footnotesize]{console}
$ nm -g librail_efr32xg24_gcc_release.a | grep -c RAIL_IEEE802154
81
$ nm -g librail_efr32xg24_gcc_release.a | grep -c RAIL_BLE
54
\end{minted}

Eighty-one 802.15.4 entry points and fifty-four Bluetooth ones, in the same binary. Nothing has to be
downloaded for the link layer of Zigbee or Thread; the question is only what the layers above the radio are
allowed to assume.

\subsection{The rules, written while Bluetooth was being brought up}

A constraint noticed later is a refactor. So the rules are stated here, at the point where the Bluetooth
path was made to work, and they bind the code that follows.

\textbf{The radio sits behind a protocol-neutral boundary.} The HAL tunes, transmits and receives. Bluetooth
and 802.15.4 are modules above it. A HAL function whose name contains \codeinline{advertising},
\codeinline{access address} or \codeinline{connection} is in the wrong layer, however convenient it looks
at the time.

\textbf{Channel numbering belongs to the protocol.} This chapter paid for that rule: RAIL's channel argument
is the \emph{physical} channel, Bluetooth's advertising channels are the pairs $(0,37)$, $(12,38)$,
$(39,39)$, and 802.15.4 numbers its channels on a different grid entirely. A caller that passes a number it
heard from a specification is passing a number into the wrong namespace.

\textbf{Services that all three protocols need are written once, and neutrally.} The timer, the persistent
store and the cipher are not Bluetooth features:

\begin{center}
\begin{tabular}{@{}lll@{}}
\toprule
Service & Needed by & Where it belongs \\
\midrule
timer over SYSRTC & every stack schedules events & the HAL \\
persistent store (nvm3) & pairing data, network keys, frame counters & one flash region, one API \\
cipher & Bluetooth CTR, 802.15.4 CCM* & a block cipher; CCM* is a mode on top \\
\bottomrule
\end{tabular}
\end{center}

The timer is the immediate consequence. The next section of work is a timer, and writing it as a
\codeinline{sl\_sleeptimer} shim would mean Zigbee needs a second driver for the same peripheral. Writing it
as a HAL timer over SYSRTC costs the same and serves three protocols.

\textbf{One state buffer is a single-protocol assumption.} \codeinline{RAIL\_AddStateBuffer} takes a buffer
per protocol instance. Code that hard-codes one global has already decided the product runs one protocol
forever.

\subsection{What the radio does not decide}

Whether Bluetooth and an 802.15.4 stack run together, time-sliced, or in separate firmware images is a
question this part answers with its 256\,kB of RAM, and the answer can wait. What cannot wait is the
boundary that keeps the question open. The test is one sentence long, and it can be asked of any new piece
of code:

\begin{quote}
\emph{Would this be identical for 802.15.4?}
\end{quote}

If yes, it belongs in the HAL or in a hardware-free library, and it must be written so that Zigbee and
Thread can use it unchanged. If no, it is a protocol module, and neither the HAL, nor the BSP, nor the other
protocol may reach into it. Applied consistently, the second radio protocol becomes an addition; applied
carelessly, it becomes a rewrite of everything this chapter built.
